dehat{A}^{-}(s,a)
        \nabla_\theta \log \pi_\theta(a \mid s)
    \right].
\end{aligned}
\label{eq:signed_actor_decomposition}
\end{equation}
For a negative sample, the update magnitude factorizes as
$\widehat{A}^{-}(s,a)
\lVert\nabla_\theta \log \pi_\theta(a \mid s)\rVert$,
separating the severity of its estimated disadvantage from its
learner-relative policy geometry.

We use \emph{coefficient magnitude} for \(\widehat A^-\), \emph{score
response} for \(\|\nabla\log\pi\|\), and \emph{update magnitude} for their
product. Section~\ref{sec:distance_strength} defines the squared remoteness
gradient as the \emph{reuse loop gain}, while \emph{aggregate force} refers to
the expectation or sum of signed updates. We use the broader term
\emph{influence} only when these distinctions are immaterial.

All aggregate quantities use the same update base measure \(\nu\). We write
\begin{equation}
    p=\mathbb E_\nu[\widehat A^+],
    \qquad
    q=\mathbb E_\nu[\widehat A^-],
    \qquad
    \rho=q/p \quad (p>0).
    \label{eq:aggregate_masses}
\end{equation}
This convention matters when comparing conditional sample fields with
population fields: sampling frequency already represented by \(\nu\) is not
applied a second time.

% ------------------------------------------------------------------
% SECTION 3: THEORETICAL FRAMEWORK
% ------------------------------------------------------------------
\section{Far-Field Dynamics of Historical Policy Updates}
\label{sec:theory}

We study how learner-relative remoteness shapes signed policy updates when
historical behavior is reused. We first characterize, in policy-output
coordinates, how the remoteness of a sample determines the strength of its
score contribution. We then show how repeated reuse turns this static
geometry into a dynamic feedback process: positive samples become
self-attenuating as the learner approaches them, whereas negative samples
become self-amplifying as the learner moves away. Finally, we analyze how
these local effects aggregate into finite stable equilibria, persistent
drift, or instability, and derive their specific manifestations in
Gaussian and categorical policies.

\subsection{Learner-Relative Remoteness and Update Strength}
\label{sec:distance_strength}

To isolate the action-side geometry, we fix a context $s$ and write
$u$ for the policy-output coordinates at that context. For a historical
action $a$, we define its learner-relative remoteness as
\begin{equation}
    D_u(a)
    =
    -\log \pi_u(a).
    \label{eq:learner_relative_remoteness}
\end{equation}
For a categorical policy, $D_u(a)$ is the surprisal of the selected
action. For a Gaussian policy, it equals a scale-dependent constant plus
one half of the squared Mahalanobis distance from the policy mean.
Thus, $D_u(a)$ measures how unlikely the historical action is under the
current learner, rather than defining a metric in the action space.

We measure the corresponding squared score norm by
\begin{equation}
    R_u(a)
    =
    \left\|
        \nabla_u \log \pi_u(a)
    \right\|_2^2.
    \label{eq:score_strength}
\end{equation}
A sample with advantage estimate $\widehat{A}(s,a)$ contributes an
output-space update of magnitude
$|\widehat{A}(s,a)|\sqrt{R_u(a)}$.
The following results characterize how learner-relative remoteness
controls this update strength.

\begin{proposition}[Distance-dependent score-function magnitude]
\label{prop:distance_strength}
Fix a context and a historical action \(a\).

For a Gaussian policy with mean \(\mu\) and fixed covariance
\(\Sigma\succ0\), let
\(C_\Sigma=\frac{d}{2}\log(2\pi)+\frac12\log|\Sigma|\) denote the
Gaussian normalization constant, which is independent of \(a\) and
\(\mu\). Then learner-relative remoteness satisfies
\begin{equation}
    D_\mu(a)-C_\Sigma
    =
    \frac12
    (a-\mu)^\top\Sigma^{-1}(a-\mu).
    \label{eq:gaussian_remoteness_mahalanobis}
\end{equation}
Moreover, the squared mean-score norm
\(R_\mu(a)=\|\nabla_\mu\log\pi_\mu(a)\|_2^2\) is bounded above and below
by positive constants times \(D_\mu(a)-C_\Sigma\), with constants
determined only by \(\Sigma\). Consequently, as standardized distance
increases, the Gaussian mean-score response grows without bound. When
\(\Sigma=\sigma^2I\), \(R_\mu(a)=2(D_\mu(a)-C_\Sigma)/\sigma^2\), so
remoteness and squared mean-score norm have the same global ordering.

For a categorical policy with logits \(z\), learner-relative remoteness is
the selected-action surprisal \(D_z(a)=-\log\pi_z(a)\). The selected-logit
score response is
\begin{equation}
    \left|
    \frac{\partial}{\partial z_a}
    \log \pi_z(a)
    \right|
    =
    1-e^{-D_z(a)}.
    \label{eq:categorical_selected_score}
\end{equation}
Thus, categorical surprisal can diverge as the selected probability
approaches zero, but the selected-logit response saturates at one. The
full logit-score norm remains bounded as well.
\end{proposition}

The proof, including the explicit Gaussian eigenvalue bounds and the full
categorical score-norm bound, is given in
Appendix~\ref{app:distance_strength}.

Proposition~\ref{prop:distance_strength} identifies the policy-family
split used below. In Gaussian policies, remoteness
corresponds to squared standardized distance and can produce unbounded
mean-score response. In categorical policies, remoteness is surprisal:
it can grow without bound as probability approaches zero, but the logit
score remains bounded.


\subsection{Self-Amplification under Historical Reuse}
\label{sec:self_amplification}

The preceding result compares samples at a fixed policy. We now consider
a fixed historical action that is repeatedly reused while the learner
changes. Write
\[
    D(u)=D_u(a),
    \qquad
    R(u)=\|\nabla_uD(u)\|_2^2,
\]
and consider the repeated single-sample update
\begin{equation}
    u_{t+1}
    =
    u_t
    +
    \eta\widehat A\nabla_u\log\pi_{u_t}(a)
    =
    u_t
    -
    \eta\widehat A\nabla_uD(u_t).
    \label{eq:single_sample_reuse_update}
\end{equation}
Here, the action \(a\) and its signed coefficient \(\widehat A\) are
held fixed over the reuse window.

\begin{theorem}[Self-attenuation and self-amplification under reuse]
\label{thm:reuse_dynamics}
Assume that \(D(u)\) is differentiable and convex on a convex
neighborhood containing the update segments, and define
\[
    D_t=D(u_t),
    \qquad
    R_t=R(u_t).
\]
If \(\widehat A<0\), then every reuse step satisfies
\begin{equation}
    D_{t+1}
    \ge
    D_t+\eta|\widehat A|R_t,
    \qquad
    R_{t+1}\ge R_t.
    \label{eq:negative_reuse_amplification}
\end{equation}
Thus, negative reuse makes the historical action increasingly remote
while preventing its score-function norm from decreasing.

If \(\widehat A>0\), and \(D\) additionally has an \(L\)-Lipschitz
gradient on the same neighborhood with \(0<\eta\widehat A\le 1/L\), then
\begin{equation}
    D_{t+1}
    \le
    D_t-\frac{\eta\widehat A}{2}R_t,
    \qquad
    R_{t+1}\le R_t.
    \label{eq:positive_reuse_attenuation}
\end{equation}
Thus, positive reuse makes the historical action increasingly compatible
with the learner while attenuating its subsequent score response.
\end{theorem}

A proof is provided in Appendix~\ref{app:reuse_dynamics}. For the two
policy-output coordinates used in the main theory, the required convexity
condition holds globally: fixed-covariance Gaussian remoteness is
quadratic in the mean, and categorical negative log-likelihood is convex
in logits.

Theorem~\ref{thm:reuse_dynamics} separates a one-time distant update from
a historical-reuse feedback loop. Under a negative coefficient, the
learner moves in the direction that increases the historical action's
remoteness; after this movement, repeated reuse preserves or increases
the sample's score response. Combined with
Proposition~\ref{prop:distance_strength}, this yields unbounded
amplification for fixed-covariance Gaussian mean updates unless the
reused action has zero mean score, and bounded but persistent
amplification for categorical logits.

\begin{proposition}[Quantitative Gaussian reuse rates]
\label{prop:reuse_rate}
For an isotropic Gaussian mean policy with covariance \(\sigma^2I\),
one fixed negative sample of mass \(q>0\), and
\(X_t=D_t-C_\Sigma>0\), the uncontrolled update satisfies
\[
 X_{t+1}=\left(1+\eta q/\sigma^2\right)^2X_t,
\]
so remoteness grows geometrically. Anticipating the taper introduced in
Section~\ref{sec:drpo}, if the same sample is weighted by
\(\exp\{-a[D_t-\tau]_+\}\), \(a>0\), then its continuous-time
remoteness is \(O(a^{-1}\log t)\), and the discrete update obeys
\[
 X_t\le a^{-1}\{\log t+\log\log t+O(1)\}.
\]
The proof is in Appendix~\ref{app:gaussian_reuse_rate}. This slowdown
does not by itself prove aggregate boundedness; that requires the
closed-loop attraction established in
Theorem~\ref{thm:gaussian_drpo_boundedness}.
\end{proposition}

\subsection{Aggregate Equilibria under Historical Reuse}
\label{sec:aggregate_equilibria}

The preceding results describe the repeated influence of an individual
historical action, but they do not imply that the aggregate policy
dynamics must diverge. Positive and negative historical samples exert
competing attractive and repulsive forces. We therefore ask whether
their aggregate effect admits a finite stable equilibrium or instead
produces persistent drift or instability.

To analyze continuous-action Gaussian and discrete-action categorical
policies within a single argument, we use a common exponential-family
output representation at a fixed context; the exact coordinates and
regularity assumptions are given in
Appendix~\ref{app:aggregate_equilibrium}. Using the aggregate masses
\(p,q,\rho\) from Equation~\eqref{eq:aggregate_masses}, let
\(\mathbf m_+\) denote the positive-only target, the mean-parameter point
selected by positive feedback alone. When \(q>0\), let \(\mathbf m_-\)
denote the magnitude-weighted moment of negative samples. It is not an
attraction target: the negative actor contribution repels the learner from
this moment.

\begin{theorem}[Aggregate attraction--repulsion regimes]
\label{thm:aggregate_equilibria}
Consider the continuous aggregate actor dynamics and the corresponding
discrete update under the regular minimal exponential-family
representation specified in
Appendix~\ref{app:aggregate_equilibrium}. Assume \(p>0\).

\textbf{Positive-only case \((\rho=0)\).}
If \(\mathbf m_+\) lies in the interior of the attainable mean-parameter
space, there is a unique finite equilibrium whose mean parameter is
\(\mathbf m_+\). It is locally asymptotically stable in continuous time
and for sufficiently small discrete step sizes.

\textbf{Positive-dominant regime \((0<\rho<1)\).}
Define the signed target
\begin{equation}
    \mathbf m^\star(\rho)
    =
    \mathbf m_+
    +
    \frac{\rho}{1-\rho}
    \left(
        \mathbf m_+-\mathbf m_-
    \right).
    \label{eq:signed_target_ratio}
\end{equation}
If \(\mathbf m^\star(\rho)\) lies in the interior of the attainable
mean-parameter space, there is a unique finite equilibrium whose mean
parameter is \(\mathbf m^\star(\rho)\). It is locally asymptotically
stable in continuous time and for sufficiently small discrete step
sizes. Holding \(\mathbf m_+\) and \(\mathbf m_-\) fixed, its displacement
from \(\mathbf m_+\) grows without bound as \(\rho\to1^-\), unless
\(\mathbf m_+=\mathbf m_-\). If the signed target leaves the attainable
mean-parameter space, no finite equilibrium exists.

\textbf{Critical regime \((\rho=1)\).}
If \(\mathbf m_+\neq\mathbf m_-\), the policy-output coordinate exhibits
persistent drift and admits no finite equilibrium. If
\(\mathbf m_+=\mathbf m_-\), every point is stationary, but no point is
an isolated asymptotically stable equilibrium.

\textbf{Negative-dominant regime \((\rho>1)\).}
Every finite stationary point, if one exists, is unstable in continuous
time and under every discrete update with positive step size. Hence,
this regime admits no locally asymptotically stable finite equilibrium.
\end{theorem}

The proof, including the exponential-family actor field and the exact
discrete-time step-size condition, is provided in
Appendix~\ref{app:aggregate_equilibrium}.
Equation~\eqref{eq:signed_target_ratio} gives the central geometric
interpretation: while positive feedback remains dominant, negative
feedback displaces the equilibrium from \(\mathbf m_+\) in the direction
\(\mathbf m_+-\mathbf m_-\), and therefore away from
\(\mathbf m_-\). For fixed moments, the multiplier
\(\rho/(1-\rho)\) increases with the relative negative strength and
diverges as \(\rho\to1^-\). These conclusions concern policy geometry
and stability; they do not imply that the resulting displacement
improves task utility.

% ------------------------------------------------------------------
% SECTION 4: METHODOLOGY
% ------------------------------------------------------------------
\subsection{Policy-Family Manifestations of Negative Dominance}
\label{sec:policy_family_manifestations}

Theorem~\ref{thm:aggregate_equilibria} establishes that negative
dominance, \(p<q\), admits no finite stable equilibrium. It does not,
however, determine how this instability manifests or whether
policy-output score norms become unbounded. This depends on the policy
family.

\begin{theorem}[Policy-family score behavior under negative dominance]
\label{thm:policy_family_manifestations}
Assume \(p<q\).

\textbf{Gaussian policy.}
Consider a Gaussian policy
\(\pi_\mu=\mathcal N(\mu,\Sigma)\) with fixed
\(\Sigma\succ0\). Its aggregate mean dynamics have a unique finite
stationary point, which is repelling. Unless initialized exactly at that
point, the distance from it diverges under both the continuous dynamics
and every discrete trajectory with fixed step size \(\alpha>0\).
Consequently, for every fixed historical action \(a\),
\begin{equation}
    \left\|
        \nabla_\mu\log\pi_\mu(a)
    \right\|_2
    =
    \left\|
        \Sigma^{-1}(a-\mu)
    \right\|_2
    \longrightarrow
    \infty.
    \label{eq:gaussian_unbounded_score}
\end{equation}

\textbf{Categorical policy.}
Consider a categorical policy
\(\pi_z=\operatorname{softmax}(z)\) with at least two actions and the
zero-sum gauge \(\mathbf 1^\top z=0\). Any finite stationary point, when
one exists, is unique and repelling. Unless initialized exactly at that
point, the continuous dynamics and every discrete trajectory with fixed
step size \(\alpha>0\) eventually leave every compact subset of the
logit space. Along some subsequence, the corresponding probability
vector approaches the boundary of the simplex. Nevertheless, for every
action \(a\) and every finite \(z\),
\begin{equation}
\begin{aligned}
    \left\|
        \nabla_z\log\pi_z(a)
    \right\|_2^2
    &=
    \left\|
        e_a-\pi_z
    \right\|_2^2 \\
    &\leq
    2
    \left(
        1-\pi_z(a)
    \right)^2
    \leq
    2.
\end{aligned}
\label{eq:categorical_bounded_score}
\end{equation}
Thus, categorical logits can become unbounded and their probabilities
can approach the support